% mage-auto-heading: foundation-model-agentic-harness-and-engineered-environment
% mage-auto-heading: probabilistic-work-needs-external-enforcement
% mage-auto-heading: enforcement-outside-the-producing-run
% Authoritative LaTeX manuscript; Markdown is a generated compatibility projection.

% <!-- part-title: The New Engineering Problem -->

% <!-- chapter-title: Agentic Machinery and Engineering Mechanisms -->

\chapter{Agentic Machinery and Engineering Mechanisms}

% <!-- section-terms: engineering-substrate, foundation-model, agentic-harness -->

% <!-- point: engineering-methods-follow-from-the-material | Engineering methods fit the substrate they act on; reinforced concrete works because concrete takes compression and steel takes tension, and agentic engineering likewise begins by asking what productive machinery entered the process and which of its properties matter. | terms: engineering-substrate -->

Engineering methods have to fit the substrate they act on. Reinforced concrete works because concrete takes compression and steel takes tension; where the reinforcement goes follows from those properties. Agentic software engineering likewise begins by asking what productive machinery has entered the engineering process and which of its properties matter.

% <!-- point: the-substrate-is-model-plus-harness-in-an-environment | The relevant substrate is not a foundation model alone but a foundation model operating through an agentic harness inside an engineered environment; the combination determines both what autonomous work can do and where engineering can act on it. | terms: engineering-substrate, foundation-model, agentic-harness -->

For MAGE, the relevant machinery is not a foundation model alone. A foundation model operates through an agentic harness inside an engineered environment. Together, they determine what autonomous work can do and where engineering can act on it.

\section{Foundation Model, Agentic Harness, and Engineered Environment}\label{heading:foundation-model-agentic-harness-and-engineered-environment}

% <!-- section-terms: engineering-substrate, foundation-model, agentic-harness, engineered-environment -->

% <!-- point: model-plus-harness-form-a-new-engineering-substrate | A foundation model supplies inexpensive reasoning; the harness supplies runtime facilities, including hooks; the engineered environment puts those facilities to use and adds project structures beyond the harness. | terms: engineering-substrate, foundation-model, agentic-harness, engineered-environment -->

% <!-- index-def: engineering-substrate -->

A \textbf{foundation model} makes broad semantic reasoning available at a marginal cost far below the skilled human labor that many such tasks previously required. Training a model can be expensive, but using an already-trained model for one additional reasoning task is comparatively cheap.

% <!-- index-def: agentic-harness -->

An \textbf{agentic harness} is the runtime that turns the foundation model’s reasoning into action. It manages the context supplied to the model, invokes tools in response to model outputs, and returns their results for further reasoning. It can also expose hooks: defined points at which additional code can run, for example before a tool executes.

% <!-- index-def: engineered-environment -->

The \textbf{engineered environment} is the project’s deliberate arrangement of representations, tools, processes, permissions, and checks around that machinery. It includes how the engineer configures and uses the harness’s facilities, as well as structures outside the harness. A hook facility belongs to the harness; installing a check through it to prevent edits to generated files is part of engineering the environment. Build checks and review procedures also belong to that environment even when they operate independently of the harness. The same harness can therefore support different engineered environments, depending on how engineers put its facilities to use.

% <!-- point: four-foundation-model-properties | Four properties of the foundation model matter to the engineering argument: broad semantic reasoning, probabilistic output, bounded working state, and inexpensive repeated reasoning. | terms: foundation-model -->

Four properties of the foundation model matter to the engineering argument:

\begin{itemize}
\tightlist
\item
  \textbf{Broad semantic reasoning.} A foundation model can interpret, synthesize, and search across open-ended domains. Give it an abstraction---a schema, state table, or policy---and it can reason about that abstraction rather than merely execute a fixed script.
\item
  \textbf{Probabilistic output.} A foundation model can resolve consequential engineering questions incorrectly, and its confidence in its own result is not independent evidence that the result is correct. The producing run cannot establish engineering assurance merely by affirming its own output.
\item
  \textbf{Bounded working state.} Active reasoning state is finite. Durable system knowledge must therefore persist outside the producing run and be supplied, retrieved, or reconstructed when later work needs it.
\item
  \textbf{Inexpensive repeated reasoning.} Search, generation, repair, and reconciliation---reasoning work that once consumed substantial engineering labor---can now be repeated at previously impractical scale.
\end{itemize}

% <!-- point: three-harness-capabilities | The harness contributes three complementary capabilities: action (tools turn outputs into operations), context management (assembling and maintaining working state), and interposition (points where the environment can observe, constrain, or deny an action). | terms: agentic-harness, injection-point -->

The harness contributes three complementary capabilities:

\begin{itemize}
\tightlist
\item
  \textbf{Action.} Tools turn model outputs into consequential operations: editing a file, running a command, calling an API, or changing some other part of the engineering system.
\item
  \textbf{Context management.} The harness assembles and maintains the working state available to the reasoner through mechanisms such as retrieval, memory, compaction, and task state.
\item
  \textbf{Interposition.} Because consequential actions pass through identifiable interfaces, the harness can expose points where the environment observes, constrains, redirects, or denies them.
\end{itemize}

% <!-- point: the-substrate-framing-is-an-emerging-consensus | None of these properties is alien to human engineering; what changes is their combination with broad inexpensive reasoning and a machine-mediated environment, and recent harness-engineering work treats autonomous capability as a property of the model–harness–environment system. | terms: engineering-substrate, agentic-harness -->

None of these properties is wholly alien to human software engineering. Humans also reason with bounded state, make uncertain judgments, act through tools, and rely on evidence outside themselves. What changes is their combination with broad, inexpensive semantic reasoning and a machine-mediated environment around its actions. Recent work on harness engineering similarly treats autonomous capability as a property of the foundation-model--harness--environment system rather than the foundation model alone.\citep{zhong-zhu2026}\citep{young-anthropic2025}\citep{tan-thinharness2026} \autoref{agentic-substrate} shows the resulting substrate. This combination changes the economics of familiar engineering problems and exposes new places to act on them.

\begin{figure}[htbp]
\centering
\includegraphics[width=\linewidth]{book/latex/assets/agentic-substrate-f28785289f.pdf}
\caption{\emph{Agentic machinery and environment.} The foundation model supplies inexpensive semantic reasoning under finite, probabilistic working conditions; the harness turns that reasoning into action, manages its context, and exposes points where the environment can observe or constrain it. These are the properties of the productive substrate the engineered environment must work with.}
\label{agentic-substrate}
\end{figure}

\section{The Reasoning Horizon}\label{heading:the-reasoning-horizon}

% <!-- section-terms: long-horizon-task, reasoning-horizon, churn, context-window -->

% <!-- point: scale-creates-the-reasoning-problem-before-an-agent-enters | Human working memory is limited even when long-term knowledge is extensive; notes, diagrams, and other external representations let engineers work with more information than they can hold actively in mind. | terms: long-horizon-task, context-window -->

Scale has always created problems. An engineer may be able to hold a small program in mind, but a large system soon exceeds that capacity. Human \textbf{working memory} is limited: we can actively hold and work with only a small amount of information at once, even when we have much more knowledge in long-term memory.\citep{cowan2001capacity} A change that looks simple in one file may require us to recall an interface elsewhere, a prior design decision, and a constraint on acceptable behavior. We use notes, diagrams, and other external representations so that we do not have to keep all of that information in our heads.\citep{norman2013}

% <!-- point: once-a-task-exceeds-active-state-question-relevant-state-must-be-reconstructed | A harness can preserve and retrieve information across invocations, extending work beyond the model's immediate context; retaining a record does not ensure that the agent will find and use it. | terms: reasoning-horizon, churn -->

% <!-- index-def: long-horizon-task -->

An agent encounters a comparable limit through its foundation model's context window: the bounded input available in a single invocation. The harness may store a plan, earlier decisions, or tool results outside that window and retrieve them later. That lets work continue across many invocations, but information needed for a particular judgment must be supplied when that judgment is made. A \emph{long-horizon task} requires a sequence of related decisions, with earlier results or commitments carried into later steps. Merely retaining a record does not ensure that the agent will find and use it.

Memory management can nevertheless expand what the whole system can work over. Research on agent memory and context management has demonstrated this through external storage, summarization, and selective retrieval.\citep{packer2023memgpt}\citep{lee2024readagent} The model's context window stays bounded, but the surrounding machinery makes more information usable across the task.

% <!-- point: availability-is-not-a-useful-representation | Available information is not always used effectively: usable context can fall short of nominal capacity, information within a window is not uniformly accessible, and repository work incurs navigation and reconstruction costs. | terms: reasoning-horizon, context-window -->

Having information available does not guarantee that it can be used effectively. Studies of long-context language models find that usable context can fall short of nominal capacity and that information within the context is not uniformly accessible.\citep{hsieh2024ruler}\citep{liu2024lostmiddle} Repository-scale studies likewise report substantial navigation and reconstruction work as relevant state spreads across a codebase.\citep{sourcegraph2026largecodebase}

% <!-- point: the-reasoning-horizon-defined | The reasoning horizon is the practical limit of work the reasoner and its supporting machinery can carry through reliably, given the available representations and resources; task demands are distinct from this system-level limit. | terms: reasoning-horizon, context-window -->

% <!-- index-def: reasoning-horizon -->

The \textbf{reasoning horizon} is the practical limit of the work a reasoner and its supporting machinery can carry through reliably, given the available representations and resources. For an agent, this depends on the foundation model together with the harness, external memory, and tools. The task imposes reasoning demands; the system has a horizon.\footnote{We borrow the term from Ramezanali and colleagues, who use it in seqBench to describe the length of a sequence a model can reliably complete.\citep{ramezanali2025seqbench} Here we extend it to the practical limits of the complete reasoning system, including its harness, external memory, and tools.}

Better machinery and better representations can extend the horizon in different ways. A harness can preserve and recover more relevant state across a task; a useful representation can reduce what must be searched, reconstructed, or inferred. Research on external representations has long shown that how information is organized affects the work needed to reason with it.\citep{larkin1987diagram}\citep{zhang1994representations} These improvements may also make work cheaper without extending the range of tasks the system can handle. \autoref{reasoning-horizon-fig} illustrates how purposeful abstraction reduces reconstruction.

\begin{figure}[htbp]
\centering
\includegraphics[width=\linewidth]{book/latex/assets/reasoning-horizon-7eba103f30.pdf}
\caption{\emph{The reasoning horizon.} At sufficient scale, a task requires more state than a finite reasoner can keep active, so some state must be reconstructed. Representation determines the cost: raw implementation carries irrelevant detail, while purposeful abstraction can keep question-relevant state smaller and cheaper to recover.}
\label{reasoning-horizon-fig}
\end{figure}

% <!-- box-family: inset -->

% <!-- inset-domain: AGENT HARNESS -->

% mage-box-prefix: Inset —

\begin{bookbox}{inset}{Context Windows and Compaction}
\label{inset-context-windows-and-compaction}
MAGE\textquotesingle s reasoning horizon appears in today\textquotesingle s agent harnesses most visibly as the problem of working within a bounded context window. The context window determines how much material the underlying model can consider directly at once, and those windows have grown substantially across generations of foundation models. Harnesses further stretch the useful working horizon by retrieving relevant state, caching or preserving selected information, and compacting an extended interaction into a smaller representation when the accumulated context becomes too large.

These mechanisms expand the amount of work an agent can carry forward, but they do so in different ways. A larger context window increases the amount of state that can remain directly available. Compaction compresses earlier state so that a longer interaction can fit within that bound, at the risk of discarding details that later matter. Retrieval and persistent state move information outside the immediate context and bring it back when needed. More generally, loop engineering can decompose work, preserve intermediate artifacts, selectively reconstruct context, and hand state between reasoning episodes. The effective working horizon of an engineered agent system can therefore be much larger than the raw context window of any single model invocation.

This is why MAGE uses reasoning horizon rather than context window. The engineering limit that matters is not a fixed token count but how much consequential reasoning state a task requires, and how effectively the surrounding system can preserve and recover that state. Today\textquotesingle s context windows are one implementation constraint on that horizon; larger models and better harnesses can move the constraint without making the underlying problem disappear.
\end{bookbox}

\section{Probabilistic Work Needs External Enforcement}\label{heading:probabilistic-work-needs-external-enforcement}

% <!-- section-terms: foundation-model, engineering-substrate -->

% <!-- point: authorship-is-separate-from-assurance | Software engineering has always separated authorship from assurance — we compile instead of trusting a prediction, test instead of accepting an explanation — because the producer's confidence is not enough. | terms: engineering-substrate, foundation-model -->

The second problem is not that agents make mistakes while human programmers do not. Software engineering has always separated authorship from assurance. We compile code instead of trusting the programmer\textquotesingle s prediction that it compiles; we run tests instead of accepting an explanation that the program should work; consequential systems add review, static analysis, formal methods, operational evidence, or independent approval because the producer\textquotesingle s confidence is not enough.

% <!-- point: enforcement-must-reside-outside-the-agent | Commodity intelligence makes that separation more important because implementation is cheap and repeatable; an agent can propose, inspect, explain, and repair, but another judgment from the same substrate does not constitute independent evidence — the environment must establish what evidence counts. | terms: engineering-substrate, foundation-model -->

Cheap, fast, repeatable agentic production makes that separation more consequential. An agent can propose a change, inspect it, explain it, and often repair it, but another probabilistic judgment from the same productive substrate does not automatically create engineering assurance. If an obligation matters, the environment must establish what evidence counts and what follows when the obligation is not met.

The U.S. defense community\textquotesingle s AI4SDLC guidance identifies the same scaling problem: as agents produce more work asynchronously, the people reviewing it can become the bottleneck.\citep{codeMilWorkflowGovernance2026} Requiring human approval is therefore only part of the answer; reviewers also need representations and evidence that let them make an informed judgment about what the machinery has produced.

Commodity intelligence is probabilistic in an engineering sense. Contemporary foundation models are learned systems whose outputs do not ordinarily carry guarantees that the engineering judgments embodied in those outputs are correct. Their present implementations may add uncertainty at several levels: generation can vary across samples; an underspecified request may admit several interpretations; intermediate reasoning and tool choices may differ; and an agentic task may compose many such judgments before producing a result. Harnesses can retry, decompose, vote, retrieve additional context, or ask another model to evaluate the first. All of these can improve reliability. None by itself turns an inferred engineering judgment into an engineering guarantee.

The details matter when evaluating a particular system, but they are too contingent to serve as the foundation of a general engineering method. Decoding can be deterministic. Errors can be correlated. One model or harness can be much better than another. Future foundation models may operate differently again. We therefore use a deliberately simple abstraction.

Suppose a realization depends on \(n\) consequential judgments, numbered 1 through \(n\). Let \(p_i\) denote the probability that judgment \(i\) is resolved acceptably. Under the simplifying assumption that these events are independent, the probability that all of them are acceptable is the product

\[
P(\text{all judgments acceptable}) = \prod_{i=1}^{n} p_i
\]

The effect is small when each judgment is highly reliable and there are only a few of them. If each judgment has a 99\% probability of being acceptable, for example, three such judgments give an overall probability of \(0.99^3 \approx 97\%\). But the chance of success falls as the number of judgments grows: \(0.99^{10} \approx 90\%\), and \(0.99^{100} \approx 37\%\). It falls faster when individual judgments are less reliable: ten judgments that are each 95\% reliable give \(0.95^{10} \approx 60\%\). Even under this deliberately favorable model, many individually reliable judgments do not compose into a comparably reliable whole.

Many techniques can improve these probabilities, most notably running agents in iterative loops that critique, test, and revise their own work. Such techniques can make individual judgments substantially more reliable. The equation is illustrative, not a model of an actual agent: real judgments are not generally independent, their probabilities may be unknown, and iterative techniques can change those probabilities. Nevertheless, the example shows why high reliability at the level of individual judgments does not establish high reliability for a system that depends on many of them.

% <!-- index-def: probabilistic-surface -->

% <!-- gloss-only: Probabilistic surface | The set of consequential decisions in a system whose correctness depends on probabilistic inference at the time they are made. A larger probabilistic surface creates more opportunities for inference error to affect the outcome, even when each individual decision is highly reliable. MAGE seeks to reduce this surface by moving suitable decisions into explicit models, deterministic mechanisms, and independently enforced obligations, while leaving genuinely judgment-dependent decisions probabilistic. -->

% <!-- point: the-probabilistic-surface-is-the-judgments-still-bought-per-instance | The consequential judgments that still depend on fallible per-instance inference form the probabilistic surface; better models, context, decomposition, harnesses, and iterative loops make those judgments more reliable but leave the obligation dependent on probabilistic judgment each time it arises. | terms: engineering-substrate, foundation-model -->

The consequential judgments that still depend on fallible per-instance inference form the system\textquotesingle s \textbf{probabilistic surface}. Better models, prompts, context, tools, decomposition, agent harnesses, and iterative loops can make those judgments more reliable. These are valuable improvements, but they leave the obligation dependent on probabilistic judgment each time it arises.

This is the immediate reliability argument for moving consequential judgments out of repeated inference, but it is not the only argument for doing so. Future agents may make many such judgments cheaply and reliably. But even granting that future, engineering disciplines require an independent basis for understanding what a system is intended to do, deciding which obligations govern it, and establishing why its behavior should be trusted. \magevalue{chapter:profession} returns to that requirement as one of engineering control and responsibility.

\section{Enforcement Outside the Producing Run}\label{heading:enforcement-outside-the-producing-run}

% <!-- point: software-engineering-can-make-an-obligation-explicit-and-enforce-it-outside-the-run | For selected obligations the environment can make the required property explicit and enforce it independently of the producing run: it can represent an architectural boundary, a dependency rule, or a safety condition and reject changes that violate it, so the judgment that established the obligation is reused across later implementation. | terms: engineering-substrate -->

Software engineering already provides mechanisms for independent enforcement. For selected obligations, the environment can make the required property explicit and enforce it outside the producing run. It can represent an architectural boundary, a dependency rule, or a safety condition, and reject changes that violate it. The judgment that established the obligation can then be reused across later acts of implementation.

% <!-- point: agentic-work-provides-identifiable-interfaces-where-enforcement-can-act | Important changes reach the system through identifiable interfaces — inputs, tool calls, edits, outputs, lifecycle transitions — where a sandbox, hook, or test can act without a rich model, and explicit representations extend the semantic reach of those same control points. | terms: agentic-harness, injection-point, engineering-substrate -->

Agentic work provides places where such enforcement can occur. Important changes reach the engineering system through identifiable interfaces: inputs, tool calls, edits, outputs, and lifecycle transitions. A sandbox can deny network access, a hook can reject a command, and a test can block a regression without requiring a rich system model. Explicit representations can extend the semantic reach of those same control points, allowing an action to be judged against architectural, behavioral, policy, ownership, or other system-level obligations.

The harness therefore supplies control points; representation determines what those control points can know. MAGE develops principled ways to use both while leaving the rest of realization open.
